% GENERATED FILE. Edit canonical lessons, then run npm run generate:books.
% canonical-source-hash: fnv1a64:22fee42bcd78613d
% canonical-lessons: TE-C231-phyanu, TE-C231-patakaru, TE-C231-niccena, TE-C231-jiguru, TE-C231-kolabadda

\chapter{Fan, Pliers, Ladder, Glue, Ruler}
\label{ch:te-fan-pliers-ladder-glue-ruler}
\hlchaptermodality{231}

\begin{chapteropening}
\textbf{By the end of this chapter:} \emph{I can say a fan, pliers, a ladder, glue and a ruler.}

Complete the last lesson of chapter 231: I can say a fan, pliers, a ladder, glue and a ruler.
\end{chapteropening}

\section[\texorpdfstring{\te{ఫ్యాను} (phyānu)}{ఫ్యాను (phyānu)}]{\te{ఫ్యాను} (phyānu) — a fan}

\label{lesson:TE-C231-phyanu}



\begin{quote}
Before the new one: say the Telugu for entertainment, then the Telugu for an actor.
\end{quote}

\subsection*{You'll want to know: \te{ఫ్యాను}}
\textbf{\te{ఫ్యాను}} — \emph{phyānu} — \textquotedblleft{}a fan\textquotedblright{}.

\begin{grammarlens}[title={What you've built}]
One more word for this chapter.
\end{grammarlens}

\subsection*{Your turn}
\begin{itemize}
\raggedright
  \item \emph{Say it:} \emph{phyānu}
  \item \emph{Say it:} \emph{phyānu}, once more
  \item \emph{Say it:} say \emph{phyānu}
  \item \emph{Recall:} say the Telugu for entertainment, then the Telugu for an actor, then say \emph{phyānu} again
\end{itemize}

\begin{tcolorbox}[breakable,colback=teal!4,colframe=teal!35!black,title={Before you move on}]
What does \te{ఫ్యాను} mean? (\textquotedblleft{}A fan\textquotedblright{}.) Say it once more.
\end{tcolorbox}
\section[\texorpdfstring{\te{పటకారు} (paṭakāru)}{పటకారు (paṭakāru)}]{\te{పటకారు} (paṭakāru) — pliers, tongs}

\label{lesson:TE-C231-patakaru}



\begin{quote}
Before the new one: say the Telugu for an actor, then the Telugu for a fan.
\end{quote}

\subsection*{You'll want to know: \te{పటకారు}}
\textbf{\te{పటకారు}} — \emph{paṭakāru} — \textquotedblleft{}pliers, tongs\textquotedblright{}.

\begin{grammarlens}[title={What you've built}]
One more word for this chapter.
\end{grammarlens}

\subsection*{Your turn}
\begin{itemize}
\raggedright
  \item \emph{Say it:} \emph{paṭakāru}
  \item \emph{Say it:} \emph{paṭakāru}, once more
  \item \emph{Say it:} say \emph{paṭakāru}
  \item \emph{Recall:} say the Telugu for an actor, then the Telugu for a fan, then say \emph{paṭakāru} again
\end{itemize}

\begin{tcolorbox}[breakable,colback=teal!4,colframe=teal!35!black,title={Before you move on}]
What does \te{పటకారు} mean? (\textquotedblleft{}Pliers, tongs\textquotedblright{}.) Say it once more.
\end{tcolorbox}
\section[\texorpdfstring{\te{నిచ్చెన} (niccena)}{నిచ్చెన (niccena)}]{\te{నిచ్చెన} (niccena) — a ladder}

\label{lesson:TE-C231-niccena}



\begin{quote}
Before the new one: say the Telugu for a fan, then the Telugu for pliers.
\end{quote}

\subsection*{You'll want to know: \te{నిచ్చెన}}
\textbf{\te{నిచ్చెన}} — \emph{niccena} — \textquotedblleft{}a ladder\textquotedblright{}.

\begin{grammarlens}[title={What you've built}]
One more word for this chapter.
\end{grammarlens}

\subsection*{Your turn}
\begin{itemize}
\raggedright
  \item \emph{Say it:} \emph{niccena}
  \item \emph{Say it:} \emph{niccena}, once more
  \item \emph{Say it:} say \emph{niccena}
  \item \emph{Recall:} say the Telugu for a fan, then the Telugu for pliers, then say \emph{niccena} again
\end{itemize}

\begin{tcolorbox}[breakable,colback=teal!4,colframe=teal!35!black,title={Before you move on}]
What does \te{నిచ్చెన} mean? (\textquotedblleft{}A ladder\textquotedblright{}.) Say it once more.
\end{tcolorbox}
\section[\texorpdfstring{\te{జిగురు} (jiguru)}{జిగురు (jiguru)}]{\te{జిగురు} (jiguru) — glue}

\label{lesson:TE-C231-jiguru}



\begin{quote}
Before the new one: say the Telugu for pliers, then the Telugu for a ladder.
\end{quote}

\subsection*{You'll want to know: \te{జిగురు}}
\textbf{\te{జిగురు}} — \emph{jiguru} — \textquotedblleft{}glue\textquotedblright{}.

\begin{grammarlens}[title={What you've built}]
One more word for this chapter.
\end{grammarlens}

\subsection*{Your turn}
\begin{itemize}
\raggedright
  \item \emph{Say it:} \emph{jiguru}
  \item \emph{Say it:} \emph{jiguru}, once more
  \item \emph{Say it:} say \emph{jiguru}
  \item \emph{Recall:} say the Telugu for pliers, then the Telugu for a ladder, then say \emph{jiguru} again
\end{itemize}

\begin{tcolorbox}[breakable,colback=teal!4,colframe=teal!35!black,title={Before you move on}]
What does \te{జిగురు} mean? (\textquotedblleft{}Glue\textquotedblright{}.) Say it once more.
\end{tcolorbox}
\section[\texorpdfstring{\te{కొలబద్ద} (kolabadda)}{కొలబద్ద (kolabadda)}]{\te{కొలబద్ద} (kolabadda) — a ruler, a measuring stick}

\label{lesson:TE-C231-kolabadda}



\begin{quote}
Before the new one: say the Telugu for a ladder, then the Telugu for glue.
\end{quote}

\subsection*{You'll want to know: \te{కొలబద్ద}}
\textbf{\te{కొలబద్ద}} — \emph{kolabadda} — \textquotedblleft{}a ruler, a measuring stick\textquotedblright{}.

\begin{grammarlens}[title={What you've built}]
One more word for this chapter.
\end{grammarlens}

\subsection*{Your turn}
\begin{itemize}
\raggedright
  \item \emph{Say it:} \emph{kolabadda}
  \item \emph{Say it:} \emph{kolabadda}, once more
  \item \emph{Say it:} say \emph{kolabadda}
  \item \emph{Recall:} say the Telugu for a ladder, then the Telugu for glue, then say \emph{kolabadda} again
\end{itemize}

\begin{tcolorbox}[breakable,colback=teal!4,colframe=teal!35!black,title={Before you move on}]
What does \te{కొలబద్ద} mean? (\textquotedblleft{}A ruler, a measuring stick\textquotedblright{}.) Say it once more.
\end{tcolorbox}
